\scanpage{69}
- \mspage{109} -
\hfill n°371
\[
(35)\qquad\langle u\lrcorner s,v\rangle=\langle s,vu\rangle
\]
quel que soit $v\in\bigwedge(M)^{*\mathrm{gr}}$.

On notera qu'il existe en général plusieurs endomorphismes distincts de $\bigwedge(M)$ ayant cette propriété, puisqu'il peut y avoir des éléments $\ne0$ de $M$ orthogonaux à $M^*$. Toutefois, si $M$ est un $A$-module \uline{libre}, il en est de même de $\bigwedge(M)$ (n°3, prop.8), donc il n'y a aucun élément $\ne0$ de $\bigwedge(M)$ orthogonal à son dual, et l'application $s\longrightarrow u\lrcorner s$ est bien la restriction à $\bigwedge(M)$ (considéré comme sous-module du dual de $\bigwedge(M)^{*\mathrm{gr}}$) de la \uline{transposée} de l'endomorphisme $v\longrightarrow vu$ de $\bigwedge(M)^{*\mathrm{gr}}$.

Nous définirons l'application linéaire $\bigwedge(M)\otimes\bigwedge(M)^{*\mathrm{gr}}\longrightarrow\bigwedge(M)$ correspondant à l'application bilinéaire $(s,u)\longrightarrow u\lrcorner s$ comme la composée
\diagram{
\begin{tikzcd}[ampersand replacement=\&,column sep=large]
\bigwedge(M)\otimes\bigwedge(M)^{*\mathrm{gr}}\arrow[r,"\Delta\otimes1"]\&\bigwedge(M)\otimes\bigwedge(M)\otimes\bigwedge(M)^{*\mathrm{gr}}\arrow[r,"1\otimes p"]\&\bigwedge(M)
\end{tikzcd}
}
où $\Delta$ a été définie au n°4 et $p:\bigwedge(M)\otimes\bigwedge(M)^{*\mathrm{gr}}\longrightarrow A$ correspond à la forme bilinéaire canonique (chap.II, 3e éd., §2, n°\quad). Si $s=x_1\wedge x_2\wedge\ldots\wedge x_p$, $u\in(\bigwedge^qM)$, il résulte déjà de cette définition que $u\lrcorner s=0$ si $q>p$ ; si $q\leq p$, la formule (9) du n°3 appliquée pour $n=2$ et avec tous les $x_{ji}$ ($1\leq i\leq n$) égaux à un même élément $x_j$ ($1\leq j\leq m$) donne, pour la composante de $\Delta(s)$ dans $(\bigwedge^{p-q}M)\otimes(\bigwedge^qM)$, l'élément $\sum_H\unsure{\rho|\varepsilon}_{K,H}x_K\otimes x_H$, où $H$ parcourt l'ensemble des parties à $q$ éléments de $\lbrack\!\lbrack1,p\rbrack\!\rbrack$ et $K$ est le complémentaire de $H$ dans $\lbrack\!\lbrack1,q\rbrack\!\rbrack$ ; d'où
\[
(36)\qquad u\lrcorner s=\sum_H\unsure{\rho|\varepsilon}_{K,H}\langle x_H,u\rangle x_K\ ;
\]
la formule (35) résulte alors aussitôt de (20). On dit que $u\lrcorner s$ est le \uline{produit intérieur gauche} de $u$ et de $s$ ; si $u\in(\bigwedge^qM)^*$, $s\longrightarrow u\lrcorner s$ est une application linéaire \uline{homogène de degré} $-q$. Lorsque $s\in\bigwedge^pM$, $u\in(\bigwedge^pM)^*$, on a $u\lrcorner s=\langle s,u\rangle$.

Pour deux éléments $u,v$ de $\bigwedge(M)^{*\mathrm{gr}}$, on a
\[
(37)\qquad(uv)\lrcorner s=u\lrcorner(v\lrcorner s)
\]
autrement dit, l'application $(u,s)\longrightarrow u\lrcorner s$ définit (avec l'addition) une structure de $\bigwedge(M)^{*\mathrm{gr}}$-\uline{module à gauche gradué} sur $\bigwedge(M)$, quand on munit $\bigwedge(M)^{*\mathrm{gr}}$ de la
\margin{68}
